% Proposition pour Emmanuel Fleury, revue après la réunion du 26 novembre 2026 et relecture du code.
% Chaque chemin cité existe dans ce dépôt : scripts/check_proposal.py le vérifie.
\documentclass[11pt,a4paper]{article}
\usepackage[T1]{fontenc}
\usepackage[utf8]{inputenc}
\usepackage[french]{babel}
\frenchsetup{AutoSpacePunctuation=false}
\usepackage[sfdefault]{plex-sans}
\usepackage[varqu,varl]{inconsolata}
\usepackage[hmargin=20mm,vmargin=15mm]{geometry}
\usepackage{booktabs}
\usepackage{array}
\usepackage{enumitem}
\usepackage{titlesec}
\usepackage{tikz}
\usetikzlibrary{arrows.meta,positioning,fit,backgrounds,calc}
\usepackage[hidelinks]{hyperref}

\setlength{\parindent}{0pt}
\setlength{\parskip}{5pt}
\setlist{itemsep=2pt,topsep=3pt,leftmargin=*}
\setlist[itemize]{label=\textbullet}
\titleformat{\section}{\large\bfseries}{\thesection.}{6pt}{}
\titlespacing*{\section}{0pt}{12pt}{4pt}
\titleformat{\subsection}{\normalsize\bfseries}{\thesubsection.}{6pt}{}
\titlespacing*{\subsection}{0pt}{8pt}{2pt}
\newcolumntype{L}[1]{>{\raggedright\arraybackslash}p{#1}}
% \path de url.sty : les _ des noms de fichiers passent tels quels.
\let\fichier\path

\begin{document}
\raggedright

{\Large\bfseries Harness d'agents et analyseurs statiques pour la détection de bugs en C}\par
\smallskip
Proposition de travail, UE Numerics. Arnaud Gomes, M2 CSI, tutorat d'Emmanuel Fleury. 27 novembre 2026,
revue après la réunion du 26 novembre et relecture du code.\par
Code : \url{https://github.com/Arnaud1404/c-bug-harness}, créé le 26 novembre 2026 pour ce travail. L'image des
analyseurs vient de c-secure-build, un projet commencé le 14 août, et garde l'historique de ses commits.

\section{Objectif}

Mesurer, à modèle fixé, combien un harness d'agents aide un petit modèle local à trouver des bugs en C. Deux
conditions sont comparées sur CASTLE : le modèle seul, avec le prompt de l'article, et le modèle avec le harness.
Les analyseurs seuls servent de référence sans modèle. Pour chaque condition je relève le temps, les tokens, le
nombre de bugs trouvés et le nombre de faux positifs.

\section{Analyseurs statiques}

Sept outils, huit exécutions, figés dans \fichier{analyzers/Dockerfile} (Ubuntu 26.10, dernières versions stables au
24 novembre 2026). Les options reprennent les wrappers de CASTLE~\cite{castle} pour chaque outil que l'article a
lancé, et \fichier{analyzers/README.md} donne la source de chaque option.

\smallskip
\begin{tabular}{@{}L{26mm}L{14mm}L{36mm}L{36mm}L{40mm}@{}}
\toprule
Outil & Version & Technique & Domaine & Faux positifs attendus \\
\midrule
Flawfinder & 2.0.20 & Cherche des noms de fonctions dangereuses & Appels comme \texttt{strcpy} ou \texttt{printf} & Élevés : ne lit pas le contexte \\
cppcheck & 2.22.0 & Suit les valeurs dans le code (data-flow analysis) & Out-of-bounds, null pointer dereferences, fuites & 9 pour 19 vrais dans CASTLE \\
GCC \texttt{-fanalyzer} & 16.2.0 & Parcourt les chemins d'exécution & Double free, use-after-free, fuites, null pointer dereferences & Ni sound ni complete, selon sa documentation \\
Clang Static Analyzer & 23.1.6 & Parcourt les chemins d'exécution & Mémoire, null pointer dereferences, division par zéro & 2 pour 13 vrais dans CASTLE \\
Infer & 1.3.0 & Trois checkers, voir~\ref{sec:infer} & Mémoire, code mort, deadlocks & Selon le checker \\
ESBMC & 8.8 & Bounded model checking & Débordements mémoire, null pointer dereferences, fuites, overflows arithmétiques & 32 pour 53 vrais dans CASTLE \\
Frama-C EVA & 33.0 & Interprétation abstraite, sound & Les erreurs à l'exécution qu'il modélise & Élevés : signale tout ce qu'il ne peut pas exclure \\
\bottomrule
\end{tabular}

\smallskip
Les chiffres de CASTLE viennent de versions plus anciennes des outils : ce sont des repères, pas des références.
GCC tourne deux fois, la seconde avec \texttt{-fanalyzer-checker=taint}, que GCC laisse désactivé par défaut.
Les commandes de compilation de CASTLE appellent \texttt{gcc} par son nom, alors que scan-build n'analyse que les
compilations qui passent par \texttt{\$CC}. Je remplace donc \texttt{gcc} par \texttt{\$CC}, sans quoi Clang
n'analyse rien.

Semgrep est retiré. CASTLE lançait Semgrep Pro en ligne avec \texttt{-{}-config auto}, alors que Semgrep CE hors
ligne ne cherchait que les patterns de mes propres règles. Ses résultats auraient mesuré mes règles plutôt que
l'outil.

\subsection{Les checkers d'Infer}\label{sec:infer}

Infer~\cite{infer} tourne avec trois checkers, chacun activé par son option et séparé au parsing
(\fichier{harness/normalize.py}). Sa documentation les donne tous les trois pour C/C++/ObjC.
\begin{itemize}
  \item \textbf{Pulse} (\texttt{-{}-pulse}) : null pointer dereferences, use-after-free, fuites
  (\texttt{MEMORY\_LEAK\_C}), valeurs non initialisées. Pulse ne signale une erreur que si elle se produit quelles
  que soient les valeurs des paramètres de la fonction. Sinon l'erreur reste latente, et elle n'est signalée que
  si un appelant remplit ses conditions~\cite{pulse}. Sur les fonctions dont il n'a pas le code, Pulse fait des
  hypothèses optimistes, ce qui peut produire des faux négatifs comme des faux positifs.
  \item \textbf{Liveness} (\texttt{-{}-liveness}) : un seul type, \texttt{DEAD\_STORE}, une valeur écrite puis
  jamais lue. Une écriture morte n'est pas un bug mais du code mort, donc je la compte à part.
  \item \textbf{Starvation} (\texttt{-{}-starvation}) : \texttt{DEADLOCK} et \texttt{STARVATION}. Il ne concerne que
  les programmes qui prennent des verrous, soit 3 des 10 programmes de CWE-362 dans CASTLE.
\end{itemize}

\section{Harness}

\begin{figure}[h]
\centering
\begin{tikzpicture}[
  font=\small,
  panel/.style={draw=black!65, line width=1.2pt, rounded corners=6pt, fill=black!5, inner sep=6pt},
  tab/.style={draw=black!65, line width=1pt, rounded corners=3pt, fill=white, font=\small\bfseries, inner sep=3pt},
  stage/.style={draw=black!70, dashed, rounded corners=4pt, fill=white, align=center, minimum height=12mm,
                text width=18.5mm, inner sep=2pt},
  result/.style={draw=black!70, rounded corners=3pt, align=center, inner sep=3pt},
  arr/.style={-{Stealth[length=2.2mm]}, line width=0.9pt},
]
\node[stage] (prog) {Programme\\CASTLE};
\node[stage, right=5mm of prog] (tools) {Analyseurs\\{\scriptsize 8 exécutions}};
\node[stage, right=11mm of tools] (triage) {\textbf{triage}\\{\scriptsize 1 tour}};
\node[stage, right=5mm of triage] (verifier) {\textbf{verifier}\\{\scriptsize 6 tours max.}};
\node[stage, right=10mm of verifier] (checker) {\textbf{checker}\\{\scriptsize 1 tour}};
\node[stage, right=10mm of checker] (writer) {\textbf{writer}\\{\scriptsize 1 tour}};
\node[result, below=6mm of writer, fill=green!12] (report) {\texttt{REPORT.md}};
\node[result, above=9mm of verifier, fill=red!10, text width=34mm] (rejected)
  {Rejetés, avec l'étape et la raison};
\node[result, below=6mm of verifier, fill=blue!8, text width=30mm, font=\scriptsize] (san)
  {\texttt{build\_san}, \texttt{run\_input}\\conteneur sans réseau};
\coordinate (top) at ($(rejected.north)+(0,5mm)$);

\draw[arr] (prog) -- (tools);
\draw[arr] (tools) -- node[below, font=\scriptsize] {findings} (triage);
\draw[arr] (triage) -- (verifier);
\draw[arr] (verifier.20) -- node[above, font=\scriptsize] {verdict} (checker.160);
\draw[arr] (checker.200) -- node[below, font=\scriptsize] {retry} (verifier.340);
\draw[arr] (checker) -- node[below, font=\scriptsize] {accept} (writer);
\draw[arr] (writer) -- (report);
\draw[arr, dashed] (triage.north) |- (rejected.west);
\draw[arr, dashed] (verifier.north) -- (rejected.south);
\draw[arr, dashed] (checker.north) |- (rejected.east);
\draw[{Stealth[length=2mm]}-{Stealth[length=2mm]}, line width=0.7pt] (verifier) -- (san);

\begin{scope}[on background layer]
  \node[panel, fit=(prog)(tools)(triage)(verifier)(checker)(writer)(report)(rejected)(san)(top)] (box) {};
\end{scope}
\node[tab] at (box.north) {Une exécution du harness sur un programme};
\end{tikzpicture}
\par\smallskip{\small Le harness tel que dessiné le 26 novembre. Au plus 2 passages verifier/checker par bug (1 \texttt{retry}).}
\end{figure}

\subsection{Un agent par tâche}

Chaque agent est un fichier Markdown du dossier \fichier{agents/}. Son en-tête fixe ses outils et son nombre
maximal de tours, et le reste du fichier est son prompt. Un agent ne voit que ce que le harness lui passe : il
ne partage ni contexte ni historique avec les autres.
\begin{itemize}
  \item \fichier{agents/triage.md} reçoit le programme numéroté et les findings des analyseurs. Il regroupe ceux qui
  désignent le même bug, garde ou rejette chaque groupe avec la ligne et la raison, et peut ajouter un bug
  qu'aucun outil n'a vu. Il n'a pas d'outil.
  \item \fichier{agents/verifier.md} reçoit un seul bug gardé. Il compile le programme sous ASan et UBSan
  (\texttt{build\_san}), puis le lance avec les arguments et l'entrée qu'il choisit (\texttt{run\_input}). Il
  répond \texttt{confirmed}, \texttt{not\_triggered} ou \texttt{refuted}. Pour les CWE sans effet visible
  à l'exécution, comme une injection SQL ou un mot de passe en dur, son prompt lui demande de le dire et de
  juger sur le code.
  \item \fichier{agents/checker.md} reçoit le bug, le verdict du verifier et les appels d'outils enregistrés
  (compilation, arguments, entrée et sorties brutes). Il ne voit pas le raisonnement du verifier. Il répond
  \texttt{accept}, \texttt{reject} ou \texttt{retry} avec ce qu'il faudrait essayer.
  \item \fichier{agents/writer.md} écrit \texttt{REPORT.md} à partir des bugs acceptés et des rejets. Il ne
  change aucun résultat, et c'est la première étape que je coupe si le temps manque.
\end{itemize}

Le checker joue le rôle d'un second avis indépendant, comme dans un LLM council. Il juge sur ce qui a réellement
tourné, pas sur ce que le verifier en dit.

\subsection{La boucle agentique}

La boucle tient dans \fichier{harness/loop.py}, environ 150 lignes sans framework. Pour un agent :
\begin{enumerate}
  \item Le harness envoie au serveur de modèle le prompt de l'agent, la tâche et la liste de ses outils.
  \item Le modèle répond par du texte ou demande un outil, avec ses arguments en JSON.
  \item Le harness lance lui-même l'outil dans l'image des analyseurs, sans réseau et avec 10 secondes de limite
  pour le programme, puis renvoie la sortie réelle au modèle. Le modèle ne lance jamais rien directement.
  \item On recommence jusqu'à une réponse sans appel d'outil. Le dernier tour n'offre plus d'outil, donc l'agent
  doit répondre.
\end{enumerate}
La boucle externe est dans \fichier{harness/pipeline.py}. Quand le checker répond \texttt{retry}, son objection
repart au verifier, qui relance sa propre boucle. Au deuxième \texttt{retry}, le bug est rejeté. Une réponse qui
n'est pas le JSON attendu compte comme vide. Chaque appel au modèle est enregistré avec ses tokens, son temps,
son motif d'arrêt, la température et les sorties des outils.

\subsection{Modèles}

Les deux modèles notés au tableau le 26 novembre : Qwen3.5 9B (quantisé Q4\_K\_M, 5,68~Go) et Gemma 4 E4B
(Q4\_0, 4,59~Go). Ils tourneront avec llama.cpp (\texttt{llama-server}, image Docker CUDA 13, build b12767) sur
une RTX 5060 Laptop de 8~Go, un modèle à la fois. Les fichiers sont figés par leur SHA-256 dans
\fichier{models.lock}. Le harness ne parle qu'à l'API compatible OpenAI du serveur, donc changer de serveur revient
à changer une URL. Gemini 4 Hydrogen ne viendra qu'ensuite, une fois mesurés les tokens qu'une exécution
consomme en local.

\section{Benchmark}\label{sec:bench}

\textbf{CASTLE}~\cite{castle} (ELTE, licence MIT), au commit \texttt{c7391c9} : 250 programmes C compilables, 10
par CWE pour 25 CWE, dont 150 vulnérables et 100 sains. La ligne du bug est dans les métadonnées, pas dans le
code. La version source de CASTLE marque cette ligne par un commentaire \texttt{\{!LINE\}}, mais la version
publiée l'enlève, comme l'en-tête du fichier.

Le score reprend le notebook de l'article (\fichier{harness/castle.py}). Un programme vulnérable est trouvé si un
finding tombe sur une ligne du bug ou porte une CWE apparentée. Il rapporte 5 points, chaque finding en trop
retire 1 point, et un programme sain sans finding rapporte 2 points. Un bonus s'ajoute selon le rang de la CWE
dans le Top 25 de MITRE. Le notebook l'accorde aussi aux programmes non trouvés dès qu'un premier vrai positif est
compté. Je garde ce calcul pour rester comparable à l'article, et je donne aussi precision et recall.

CrashBench et SWE-bench, envisagés dans la version précédente, sont écartés : le schéma du 26 novembre ne garde que CASTLE.

\section{Protocole}\label{sec:proto}

\begin{itemize}
  \item \emph{Analyseurs} : l'union des huit exécutions, une ligne comptée une fois, puis chaque outil et chaque
  checker d'Infer seuls.
  \item \emph{LLM seul} : le prompt et le format de réponse de l'article CASTLE, sans outil.
  \item \emph{LLM + harness} : le harness complet. Le score est calculé après le triage, après le verifier et
  après le checker, ce qui donne l'apport de chaque agent avec une seule exécution.
\end{itemize}

\begin{itemize}
  \item Mesures : pour chaque condition et chaque modèle, le temps total, les tokens de prompt et de réponse
  relevés par \texttt{llama-server}, les réponses invalides, le nombre de bugs trouvés et de faux positifs, le
  score CASTLE, la precision et le recall avec un intervalle de Wilson à 95\,\%. Un test de McNemar exact
  compare les deux conditions avec modèle.
  \item Ordre : d'abord un pilote stratifié de 25 programmes (un par CWE, 15 vulnérables et 10 sains) pour mesurer
  temps, tokens et comparer température 0 aux réglages par défaut de \texttt{llama-server}, puis les 250 si le
  temps le permet.
  \item Contamination : CASTLE est public, donc un modèle a pu le voir à l'entraînement. Ce biais touche les
  conditions avec modèle mais pas les analyseurs.
  \item Variance : sur les 25 programmes du pilote, une seconde exécution aux réglages par défaut mesure l'écart
  entre deux passages, face au passage déterministe à température 0.
\end{itemize}

\section{État au 27 novembre et calendrier}

Ce qui existe : l'image des analyseurs et son test (\fichier{analyzers/test.sh}, les huit exécutions finissent),
la normalisation de leurs rapports (\fichier{harness/normalize.py}), le score CASTLE, la boucle, les deux outils,
les quatre agents et le pipeline, relus ligne par ligne le 27 novembre. Les tests de \fichier{tests/test_harness.py}
tournent sans GPU : la boucle face à un serveur qui rejoue des réponses écrites à l'avance, et le pipeline avec ses
allers-retours. Les deux outils ont tourné une fois sur un vrai programme de CASTLE. Aucun modèle n'a encore tourné,
et \fichier{scripts/serve.sh} n'est pas encore testé.

Début décembre : premier passage sur 25 programmes et mesure des tokens. Mi-décembre : les 250 programmes.
Début janvier : tableaux et rédaction. Rendu visé fin janvier 2027 : un rapport d'au moins une douzaine de pages,
le code, les prompts, les sorties brutes et les scripts qui recalculent les tableaux.

Limites : 250 petits programmes ne donnent pas un taux de détection sur du code réel. Les résultats ne valent que
pour ces deux modèles, à cette quantisation. Un modèle de 9 milliards de paramètres peut mal suivre un format
d'appel d'outil, ce qui se verra dans la part de réponses invalides.

{\small
\begin{thebibliography}{9}\setlength{\itemsep}{0pt}
\bibitem{castle} R.~A.~Dubniczky, K.~Z.~Horvát, T.~Bisztray et al.
\emph{CASTLE: Benchmarking Dataset for Static Code Analyzers and LLMs towards CWE Detection}. arXiv:2503.09433,
2025. \url{https://github.com/CASTLE-Benchmark/CASTLE-Benchmark}
\bibitem{infer} Meta. \emph{Infer 1.3.0, documentation des checkers}. \url{https://fbinfer.com/docs/all-checkers}
\bibitem{pulse} Q.~L.~Le, A.~Raad, J.~Villard, J.~Berdine, D.~Dreyer, P.~W.~O'Hearn. \emph{Finding Real Bugs in
Big Programs with Incorrectness Logic}. OOPSLA, 2022.
\end{thebibliography}}

\end{document}
